Este apêndice complementa a justificativa resumida apresentada no
Capítulo~\ref{chapter:desenvolvimento}, registrando os critérios que levaram à
adoção da família ESP32 e do ESP-NOW no WCAN. A discussão é mantida fora do corpo
principal porque detalha alternativas de plataforma, alimentação e comunicação
sem fio que orientaram o projeto, mas não constituem a contribuição central do
trabalho.

\section{Microcontrolador}

A escolha do \emph{hardware} partiu das restrições impostas ao nó sensor. O conjunto de processamento, rádio e alimentação precisava ocupar uma área de poucos centímetros quadrados para viabilizar a fixação em locais confinados, como o interior da roda ou o perfil interno de asas aerodinâmicas, sem interferir na montagem ou na aerodinâmica. Essa restrição favorecia placas integradas, nas quais microcontrolador, rádio e circuito de bateria ficam no mesmo módulo.

O segundo critério foi o custo unitário. Como a proposta depende de vários nós de comunicação, o orçamento de uma equipe de Fórmula SAE exigia uma solução da ordem de dezenas a pouco mais de cem reais por nó em protótipo, sem incluir sensores específicos. A estimativa da Tabela~\ref{tab:bom-no-sensor}, com valores levantados em maio de 2026, mostra a base dessa ordem de grandeza e explica por que módulos comerciais ou industriais, normalmente muito mais caros, não eram uma opção natural.

A terceira restrição era operacional. A equipe já desenvolve \emph{firmware} embarcado em C/C++ para STM32, principalmente com \emph{STM32CubeIDE} e \emph{STM32CubeMX}, usados nos módulos atuais do sistema de aquisição. Por isso, uma plataforma nova só seria viável se mantivesse um fluxo de desenvolvimento próximo à prática já dominada pela equipe.

\begin{table}[htbp]
\centering
\small
\setlength{\tabcolsep}{5pt}
\renewcommand{\arraystretch}{1.12}
\begin{tabular}{p{0.38\linewidth}|p{0.22\linewidth}|p{0.30\linewidth}}
\hline
Item do nó de comunicação & Estimativa & Observação \\
\hline
Placa ESP32-C3 compacta & R\$40--70 & Integra microcontrolador, rádio e carregamento de bateria. \\
\hline
Conectores, cabos, placa auxiliar e passivos & R\$15--30 & Varia conforme sensor e fixação. \\
\hline
Bateria reaproveitada ou célula pequena & R\$0--25 & Valor de protótipo; exige célula aprovada pelos critérios de inspeção. \\
\hline
Invólucro e fixação simples & R\$10--25 & Não inclui encapsulamento automotivo final. \\
\hline
Total aproximado & R\$65--150 & Estimativa de bancada, sem sensor específico. \\
\hline
\end{tabular}
\caption{Estimativa de custo do nó de comunicação usado como referência de protótipo, em maio de 2026.}
\label{tab:bom-no-sensor}
\end{table}

Nesse contexto, a primeira opção analisada foi manter a família STM32, já utilizada no restante do sistema de aquisição da equipe. Embora a linha STM32WB suporte \emph{Bluetooth Low Energy} e 802.15.4, ela não possui Wi-Fi nativo \cite{st_stm32wb_series}. Uma alternativa seria combinar um microcontrolador STM32 convencional com um radiotransmissor Wi-Fi externo. Contudo, essa solução exigiria separar processamento, rádio e circuito de bateria em blocos distintos, o que aumentaria significativamente o \emph{footprint} e a complexidade da placa. Logo, a restrição por manter o nó de comunicação em formato compacto e direto tornou essa via desfavorável, motivando a busca por um componente com todos esses subsistemas altamente integrados.

Diante desses critérios, a família ESP32, da Espressif, foi adotada por reunir
rádio Wi-Fi integrado, processamento embarcado, placas compactas de baixo custo e
desenvolvimento de \emph{firmware} em C/C++. Essa
combinação atende simultaneamente às restrições físicas, à faixa de custo
estimada e à necessidade de manter um fluxo de desenvolvimento próximo ao já
utilizado pela equipe. O Beetle ESP32-C3, por exemplo, mede
25 por 20,5 mm, inclui gerenciamento de bateria na própria placa \cite{dfrobot_beetle_esp32c3} e tem
massa de 4 g, aferida em bancada pelo autor. Essas características tornam a
plataforma compatível com sensores instalados em rodas, componentes aerodinâmicos
ou pontos de difícil acesso.

Além da escolha do microcontrolador, a alimentação dos nós foi guiada por custo e 
facilidade de acesso. Aproveitando o circuito de gerenciamento de bateria do Beetle, 
a equipe optou por usar baterias de lítio recarregáveis, que oferecem boa densidade 
energética e são amplamente disponíveis.

As opções de mercado, porém, não apresentaram um bom equilíbrio entre custo, tamanho
e capacidade. Por isso, a equipe escolheu células reaproveitadas
para o protótipo final. A motivação é dupla: reduzir o custo do nó sensor e dar uso técnico a baterias
que ainda estejam em boas condições, evitando o descarte prematuro de células
perfeitamente utilizáveis. Para nós em que massa e volume não são restritivos,
ou em sensores com maior demanda de energia, foram consideradas células
reaproveitadas de baterias de notebook, por oferecerem maior capacidade. Para nós
em que massa e volume importam, como os instalados em regiões mais compactas,
foram consideradas células de cigarros eletrônicos descartados, por serem menores
e mais leves. Essa possibilidade foi considerada inicialmente a partir de uma série de vídeos de
demonstração sobre reaproveitamento de células de cigarros eletrônicos
descartados, citada aqui apenas como registro da inspiração prática, e não como
referência técnica de segurança ou validação elétrica.\footnote{Série de
demonstrações: Chris Doel, vídeos sobre reaproveitamento de células de cigarros
eletrônicos descartados, YouTube, 2024. Disponíveis em: \url{https://www.youtube.com/watch?v=HwoZg3BCigU}; \url{https://www.youtube.com/watch?v=dy-wFixuRVU}; \url{https://www.youtube.com/watch?v=kMiJdfgIfqI}; \url{https://www.youtube.com/watch?v=VcVp9T8f_W4}; \url{https://www.youtube.com/watch?v=ehp23hrrEHY}. Acesso em: 15/05/2026.}

O uso de células reaproveitadas, entretanto, exige critérios de aceitação mais rígidos
do que a compra de uma célula nova e rastreável. Para uma versão operacional, cada
célula deve ter procedência conhecida sempre que possível, tensão inicial dentro da
faixa esperada, corrente de descarga compatível com o consumo do nó, proteção contra
sobrecarga, sobredescarga e curto-circuito, carregamento por circuito adequado para
célula única, fixação mecânica que evite perfuração ou esmagamento e invólucro que
impeça contato acidental com terminais. Células com inchaço, dano mecânico,
aquecimento anormal, tensão fora da faixa esperada, histórico incompatível com o uso
pretendido ou comportamento instável em ensaio de carga devem ser descartadas. Essa
cautela é compatível com a orientação de que baterias removidas de uma aplicação só
devem ser reutilizadas quando puderem ser gerenciadas com procedimentos de segurança
e rastreio, e que baterias danificadas ou inseguras não devem ser reutilizadas
\cite{epa_lithium_battery_recycling_faq}.

\section{Protocolo de Comunicação Sem Fio}

Foram considerados Bluetooth Low Energy, Wi-Fi convencional e ESP-NOW. O BLE foi
descartado porque exige descoberta e estado de conexão antes da troca de dados, o
que contraria o requisito de transmissão imediata. O Wi-Fi convencional oferece boa
vazão, mas depende de associação a uma rede, configuração de infraestrutura,
endereçamento IP e estabelecimento de sessão. Comparações experimentais também
mostram que alcance, latência, consumo e vazão variam bastante entre ESP-NOW,
Wi-Fi e Bluetooth, tornando a escolha dependente do uso \cite{eridani2021_espnow_wifi_bluetooth}.

O ESP-NOW foi escolhido por permitir comunicação direta entre dispositivos ESP32 sem
ponto de acesso nem endereço IP. Ele opera sobre a camada física do Wi-Fi e usa
endereçamento por MAC. O limite de \emph{payload} depende da versão disponível na
\emph{build}: a versão clássica do ESP-NOW trabalha com quadros menores, enquanto
o ESP-NOW v2, usado na campanha final, amplia esse espaço. Por isso, no WCAN o
tamanho máximo útil do pacote é derivado em tempo de compilação, e não fixado como
constante do projeto \cite{espressif_espnow_programming_guide,pasic2021_espnow_esp32}.
Essa combinação atende melhor aos requisitos de início rápido, baixo \emph{overhead} e
operação sem infraestrutura.

Assim como em outros protocolos sem fio, o ESP-NOW trabalha com três modos de comunicação:
\begin{itemize}
    \item \textbf{\emph{Unicast}}: comunicação ponto a ponto entre um transmissor e um receptor específico, 
    usando o endereço MAC do receptor e com status de envio no nível MAC.
    \item \textbf{\emph{Multicast}}: comunicação de um transmissor para um grupo de receptores inscritos, 
    com status de envio no nível MAC.
    \item \textbf{\emph{Broadcast}}: comunicação de um transmissor para todos os dispositivos no alcance, 
    sem confirmação de recepção.
\end{itemize}

A escolha do meio de transporte não é direta, pois o sistema precisa conciliar entrada dinâmica de nós, múltiplos receptores e confirmação de entrega. A opção mais natural é o \emph{broadcast} com ACK na aplicação, por manter os dados disponíveis a qualquer nó interessado, de forma próxima ao comportamento lógico do CAN. O \emph{unicast}, por outro lado, foi descartado porque só atenderia múltiplos receptores repetindo a mesma mensagem para cada endereço MAC conhecido, o que aumentaria o uso do rádio e reproduziria em software a lógica do \emph{multicast}. O envio direto por grupo também não resolve sozinho a entrada dinâmica de nós; por isso, a alternativa baseada em grupo foi implementada de forma híbrida, iniciando em \emph{broadcast} e, à medida que os dispositivos trocam informações, descobrindo os endereços MAC dos receptores e migrando as transmissões seguintes para o grupo cadastrado. Assim, a avaliação prática comparou dois sistemas de transporte: primeiro, o \emph{broadcast} com ACK na aplicação; segundo, o \emph{multicast} híbrido. Essa investigação justifica-se porque o projeto foca na expansão da aquisição de dados, e não na substituição da rede CAN para controles críticos \cite{fernandes2022_formula_sae_wifi_can}.